\section{Introduction}
\label{sec:introduction}

\subsection{Short-Burst as a Meta-Algorithm}

Short-Burst~\citep{cannon2022,g6paper} is a meta-algorithm for compactness
optimization in redistricting.
It runs $n_{\mathrm{bursts}}$ short Markov chains of length
$\ell_{\mathrm{burst}}$, recording only the endpoint of each chain, then
selects the endpoint with the best (lowest) edge-cut.
The inner chain restarts from the previous winner after each burst, so
the trajectory progresses by finding successively better plans---a form
of iterated local search embedded in a Markov chain.

The critical observation is that \emph{the inner chain type is a parameter}.
\citet{cannon2022} demonstrate \sburst{} with a \recom{} inner chain, but
there is nothing in the Short-Burst framework that requires \recom{}.
Any ergodic Markov chain on the space of valid redistricting plans can
serve as the inner chain.
Swapping the inner chain type changes the distributional properties of the
endpoint without altering the burst-selection logic.

\subsection{Why Chain Type Matters}

The standard \recom{} chain~\citep{deford2021} is not reversible with
respect to any simple target distribution.
It proposes new plans at high acceptance rates ($60$--$80\%$) but does not
correct for the asymmetry in proposal probabilities.
For pure compactness optimization this is irrelevant---we want good plans,
not samples from a specific distribution.
But for ensemble sampling, the lack of distributional correctness means
the endpoint distribution of a \recom{}-based \sburst{} chain is biased
toward whatever plans \recom{} tends to overproduce.

ForestRecom~\citep{mccartan2020,g9paper} and MergeSplit~\citep{janson2022,g10paper}
both apply a Metropolis--Hastings correction to the spanning-tree proposal,
making their stationary distributions approximately uniform over valid plans
(subject to the population-balance constraint).
When used as \sburst{} inner chains, this property carries over to the
burst endpoints: the selected endpoint is not merely the best plan the
(biased) chain encountered, but the best plan from a set of
approximately-correctly-distributed endpoints.

This paper studies what happens when we compose \sburst{} with \frecom{}
and \ms{}: we obtain \sbf{} and \sbms{}.

\subsection{Three Contributions}

\begin{enumerate}
  \item \textbf{Implementations.}
        We give the complete specifications of \sbf{} and \sbms{} as
        implemented in \texttt{redist-cli}, including the deterministic
        seed schedule with chain-type-specific prefixes and the two-RNG
        step protocol required by ForestRecom and MergeSplit
        (Section~\ref{sec:algorithm}).

  \item \textbf{Comparative evaluation.}
        We compare standard \sburst{} (G.6), \sbf{}, and \sbms{} on
        North Carolina ($k=14$), Wisconsin ($k=8$), and Texas ($k=38$)
        at equal compute budget (1000 total steps).
        The key finding is that at equal budget, standard \sburst{}
        achieves comparable or slightly better final edge-cut because its
        higher acceptance rate explores more plans per burst
        (Section~\ref{sec:comparison}).

  \item \textbf{Per-burst acceptance rate diagnostic.}
        We introduce $\alpha_{\mathrm{burst}}$, the fraction of burst
        steps that are accepted within a given burst.
        This diagnostic is recorded in the audit manifest and quantifies
        the effective burst length for each chain type.
        Lower $\alpha_{\mathrm{burst}}$ means fewer accepted plans per
        burst, so the nominal burst length $\ell_{\mathrm{burst}}$
        overstates the effective exploration per burst
        (Section~\ref{sec:algorithm}).
\end{enumerate}

\subsection{Paper Structure}

Section~\ref{sec:background} reviews the three chain types---\sburst{},
\frecom{}, and \ms{}---and the seed architecture.
Section~\ref{sec:algorithm} gives the formal pseudocode for \sbf{} and
\sbms{} and defines the per-burst acceptance rate diagnostic.
Section~\ref{sec:comparison} reports the empirical comparison on NC, WI,
and TX and explains the acceptance-rate mechanism behind the results.
Section~\ref{sec:when-to-use} gives usage guidance for practitioners.
Section~\ref{sec:conclusion} concludes.
